\subsection{Longitudinal Administrative Data in Healthcare}
\label{sec:longitudinal_data}

Longitudinal administrative data in healthcare are repeated observations of the same patients over time. Usually they are stored in Electronic Health Records (EHRs). EHRs contain a lot of clinical and administrative information. For example diagnoses, medical procedures, prescribed medications, and laboratory results~\cite{carrasco2023prediction}. Cross-sectional data show only a snapshot at one moment. Longitudinal records show how health conditions develop and in which order medical interventions happen~\cite{moglia2025artificial}.

Temporal trends are the main value of longitudinal EHR data. Single-time-point measurements usually have much less predictive power. When it is known how a patient's condition evolves, it is possible to study disease trajectories and also dynamics between sequential clinical events~\cite{carrasco2023prediction}. So long-term EHR histories are used more and more in epidemiological research and predictive modeling. For example to estimate risk of disease onset, to predict all-cause mortality, or to evaluate healthcare quality measures like hospital readmissions~\cite{cosco2017operationalising, carrasco2023prediction}.

To standardize all this administrative data, healthcare systems use structured medical coding frameworks. Globally, diagnoses are classified using the International Classification of Diseases (ICD-10). In the Czech Republic, healthcare providers report hospitalizations using CZ-DRG (Czech Diagnosis-Related Group) classification, which is managed by the Institute of Health Information and Statistics (ÚZIS)~\cite{uzis_czdrg_2022}. Individual medical interventions are reported to health insurance providers using the national ``List of Health Procedures with Point Values'' (\textit{Seznam zdravotních výkonů s~bodovými hodnotami})~\cite{vzp_coding_2024}. For machine learning, these chronological sequences of standardized codes are used as input features.

But longitudinal EHR datasets also have several problems. Data points are generated only when a patient comes to the healthcare system, for example during a clinic visit or hospitalization. So records are collected at irregular time intervals~\cite{moglia2025artificial}. Another problem is that the data are usually sparse, high-dimensional, and heterogeneous. They often contain missing values across different types of clinical variables~\cite{cascarano2025integrating}. Methods for sequence modeling are often taken from natural language processing (NLP) and bioinformatics. But EHR data need specialized architectures and preprocessing, because they have specific clinical properties and the structure is irregular~\cite{carrasco2023prediction}.

\subsection{Machine Learning for Sequential Health Data}
\label{sec:ml_sequential_ehr}

Machine learning on longitudinal EHR data is usually framed as a prediction problem. The patient's history is a time-ordered sequence of clinical events. It is used to estimate risk of one or more outcomes, for example a future diagnosis, adverse event, mortality, or healthcare utilization~\cite{carrasco_ribelles_2023_jamia, moglia_2025_bmc}.

\subsubsection{Problem formulation and data representations}
\label{subsec:problem_formulation}

It is necessary to decide how to represent a patient trajectory. Many approaches build sequences of visits (or time bins). Each element contains sets of medical codes (diagnoses, procedures, medications) and sometimes also continuous measurements. The chronological order is kept. This representation is discussed in the literature as a practical compromise between clinical fidelity and model tractability for longitudinal EHR prediction tasks~\cite{carrasco_ribelles_2023_jamia, moglia_2025_bmc}. In administrative settings, the longitudinal sequence is usually a sparse, high-dimensional stream of coded events. It is not a dense physiological time series. This ``code sequence'' view is close to methods originally developed for other sequential data, for example language modeling~\cite{pubmed_37652934}.

\subsubsection{Model families for longitudinal EHR sequences}
\label{subsec:model_families}

The most frequently reported deep learning architectures for longitudinal EHR prediction are recurrent models (RNNs) like LSTMs and GRUs, convolutional models adapted to sequential inputs, and transformer-based models with self-attention. These families appear repeatedly in systematic and scoping reviews of AI prediction on longitudinal EHR data~\cite{carrasco_ribelles_2023_jamia, moglia_2025_bmc}. Transformers are attractive for longitudinal administrative sequences because self-attention can model long-range dependencies between events. For example, transformer frameworks were evaluated on claims-style pediatric longitudinal data for downstream prediction tasks~\cite{nature_claims_transformer_2022}. Classical probabilistic models and bag-of-codes models are still widely used, especially when interpretability and limited training data are the main constraints. These are described in the next two subsections, because they are the families used in this work.

\subsubsection{Hidden Markov Models for Longitudinal Health Data}
\label{subsec:hmm_ehr}

Hidden Markov Models (HMMs) are one of the basic probabilistic frameworks for modeling sequential processes with latent structure. In an HMM the system is in one of a discrete set of unobserved (hidden) states at each time step. Observable outputs are generated according to emission probabilities given the current state. State transitions are given by a Markov transition matrix~\cite{rabiner1989hmm}. This formulation is suitable for clinical longitudinal data, because the underlying disease stage of a patient is not directly observed. It has to be inferred from the sequence of recorded clinical events.

HMMs were used in healthcare to model disease progression for several conditions. Gupta et al.\ applied HMMs to continuous mortality risk monitoring in ICU patients. Their temporal model was able to track health state transitions. It reached AUROC of 0.87, which was much better than non-temporal classifiers like logistic regression and random forest under the same evaluation~\cite{gupta2020hmm}. Verma et al.\ worked with the specific problem of administrative EHR data, where patients are observed only at irregular intervals, when they come to the healthcare system. They used a continuous-time HMM on a longitudinal cohort of 76,888 patients with COPD. They showed that such models can give interpretable representations of disease progression dynamics from administrative claims. It is directly useful for hypothesis generation in epidemiological research~\cite{verma2018copd}. More recently, Das modeled progression of chronic diseases (diabetes, CKD, and cardiovascular disease) using HMMs trained on longitudinal EHRs. The model was able to infer latent disease stages and transition probabilities between them~\cite{das2025hmm}.

HMM outputs are also interpretable, which is another advantage. The transition matrix gives an explicit probabilistic description of disease dynamics. The sequence of inferred hidden states can be examined to describe individual patient trajectories. Kwon et al.\ developed DPVis, a visual analytics system which exposes HMM parameters and state sequences as interactive visualizations. Clinical scientists can explore disease progression pathways across chronic diseases like type 1 diabetes, Parkinson's disease, and COPD~\cite{kwon2020dpvis}. It means HMM outputs can be inspected directly by domain experts. Internal representations of deep learning models usually cannot. It fits user-centered design goals of this work.

A scoping review of analytical methods for identifying care sequences from health utilization data confirms that Markov models, including HMMs, are one of the most frequently used families of methods for sequence analysis in administrative healthcare data (20\% of reviewed studies). The other frequent families are clustering and pattern mining~\cite{flothow2023scoping}. HMMs work well on sparse, irregularly sampled event sequences. Also, their output is interpretable. So for predictive modeling on administrative procedure code sequences (like the ones used in this work) HMMs are a practical and well-validated choice.

\subsubsection{Bag-of-Codes Representations and Classical Classifiers}
\label{subsec:bag_of_codes}

Transformer-based and recurrent architectures get a lot of attention for sequential EHR modeling. But classical machine learning methods on aggregate, order-agnostic representations of patient histories are still competitive. They are often preferred when interpretability, computational efficiency, or limited training data are the main constraints~\cite{humbert2021tfidf}.

A common approach is to represent each patient's longitudinal history as a high-dimensional binary or frequency-weighted vector over the vocabulary of medical codes observed during a defined time window. It is similar to the bag-of-words (BoW) model in natural language processing. The sequential order of codes is dropped. The patient is represented by aggregate presence or weighted frequency of codes across all encounters in the window. When Term Frequency-Inverse Document Frequency (TF-IDF) is used for weighting, codes that appear frequently in one patient's record but rarely across the population get higher weight. So clinically more informative codes are emphasized over very common administrative codes~\cite{humbert2021tfidf, chen2023scope}.

These representations were used successfully with both logistic regression and random forest on administrative data. Mukherjee et al.\ (the SCOPE study) used a cohort of more than 500,000 patients. They trained a calibrated random forest on aggregated past diagnosis codes. It reached median AUROC of 0.904 across 583 ICD-10 prediction targets. It outperformed a recurrent neural network baseline on overall discrimination. It was also much more interpretable through SHAP feature attributions~\cite{chen2023scope}. Similarly, Kalloo et al.\ used high-dimensional administrative claims data, containing thousands of ICD and procedure codes. They showed that standard machine learning approaches can predict early hospital readmission with acceptable discrimination on such data~\cite{kalloo2014claims}. So the code vocabulary together with well-regularized linear or ensemble classifiers is often enough to capture clinically relevant patterns in administrative sequences. Deep sequence models are not always needed.

These classical pipelines have another advantage. They are compatible with post-hoc interpretability techniques. Logistic regression coefficients directly show which codes most strongly predict each outcome. Random forest feature importances and SHAP values give patient-level attribution of a prediction to specific historical codes. It is exactly the type of output which the web application described in this work is designed to visualize for non-technical clinical users.

\subsubsection{Transformer-based foundation models for EHR data}
\label{subsec:foundation_models_ehr}

The success of BERT in natural language processing led to several domain-adapted foundation models for clinical data. ClinicalBERT was pretrained on free-text clinical notes from MIMIC-III and fine-tuned for 30-day readmission prediction~\cite{huang2019clinicalbert}. BEHRT adapted the BERT pretraining framework directly to structured EHR code sequences. Each token was a diagnosis code and the input embedding also included age, visit segment, and positional encoding. BEHRT was pretrained on 1.6 million UK primary care patients. It reached 8--13\% improvement in average precision over previous deep learning baselines like RETAIN and Deepr~\cite{li2020behrt}. Med-BERT used a similar strategy but at much larger scale, pretrained on 28 million US patients with a vocabulary of over 82,000 ICD-9 and ICD-10 codes. It improved AUROC by 1--6\% over non-pretrained baselines, and the benefit was biggest when the fine-tuning dataset was small~\cite{rasmy2021medbert}. A direct comparison on 6.3 million UK primary care patients confirmed that transformer representations outperform bag-of-words baselines, but the absolute improvement over well-regularized binary disease indicator vectors was small~\cite{beaney2024comparing}. A systematic review of 84 clinical foundation models showed that most are pretrained on narrow datasets (usually MIMIC-III), evaluated on tasks that do not reflect real deployment, and rarely validated externally~\cite{wornow2023shaky}. So transformer-based foundation models are the current frontier for EHR representation learning, but their direct application to small-scale, non-English administrative datasets like Czech DRG or the national procedure list (\textit{Seznam zdravotních výkonů}) needs careful adaptation. Classical models like HMMs and TF-IDF classifiers are still competitive when interpretability, data volume, and coding-system specificity are the main constraints, which was the situation in this work.

\subsubsection{Key challenges in EHR sequence modeling}
\label{subsec:ehr_challenges}

Several properties of EHR data make sequence modeling harder than for standard time series. Clinical events are recorded only when a patient interacts with the healthcare system, so the temporal spacing between observations varies a lot across patients and during follow-up of one patient. It makes direct application of standard recurrent architectures complicated, because they implicitly assume uniform time steps. Time-aware model variants or explicit time-encoding strategies are used to handle this~\cite{carrasco_ribelles_2023_jamia, moglia_2025_bmc}. Administrative EHR sequences are also represented over large vocabularies of standardized codes (ICD-10 diagnoses, DRG groups, procedure codes). So the input matrices are high-dimensional and sparse. For most visits the majority of possible codes are absent. Embedding layers and pretraining strategies are used to compress sparse code spaces into dense, semantically meaningful representations~\cite{nature_claims_transformer_2022, pubmed_37652934}.

Missingness in EHR data is rarely random. A patient may not have a recorded diagnosis not because the condition is absent, but because it was never tested, or because it was managed in another institution. In administrative datasets the data-generating process is driven by reimbursement rules, not by clinical completeness. So there are systematic biases that differ from research-grade cohorts. This so-called ``informative missingness'' can make predictive models learn spurious associations between coding intensity and the target outcome~\cite{carrasco_ribelles_2023_jamia, moglia_2025_bmc}. Another frequent problem is temporal leakage. It happens when features derived from events that occurred after the prediction time point accidentally influence training. It can produce artificially high performance estimates which do not generalize to prospective clinical use. Systematic reviews of AI models on longitudinal EHR data reported both label inconsistency and inadequate temporal separation as widespread methodological problems~\cite{carrasco_ribelles_2023_jamia, moglia_2025_bmc}.

\subsubsection{Evaluation, validation, and reporting practices}
\label{subsec:eval_reporting}

Performance of machine learning models on longitudinal EHR data is usually evaluated in two main dimensions: discrimination and calibration. Discrimination metrics like AUROC and AUPRC measure how well the model can distinguish positive from negative cases. Calibration measures agreement between predicted probabilities and actual observed frequencies. A highly discriminative model can still be poorly calibrated, and the predicted probabilities can then be misleading for clinical decision-making~\cite{calib_discrim_ecg_2023, carrasco_ribelles_2023_jamia}. There are a lot of deep learning models for EHR sequences, but the methodological quality of published studies is still a problem. Systematic reviews based on the TRIPOD and TRIPOD-AI reporting guidelines found that AI models trained on EHRs often do not describe data preprocessing, handling of missing data, prediction window definition, and code sharing well enough. Most models are also validated only internally, with simple train-test splits. External validation across different healthcare institutions or patient cohorts is rare, so geographical and temporal dataset shift is not assessed~\cite{carrasco_ribelles_2023_jamia, bmj_tripod_ai_2024}. Another problem is that there are not many standardized longitudinal benchmark datasets. MIMIC-III and MIMIC-IV are widely used but are restricted almost only to ICU and emergency department visits. They represent very acute, short-term physiological trajectories and not long-term administrative histories. Recent initiatives like the EHRSHOT benchmark try to provide standardized, few-shot evaluation frameworks designed for longitudinal, non-ICU administrative data~\cite{wori_ehrshot_2023}.

\subsubsection{Interpretability and communicating model outputs}
\label{subsec:interpretability}

For clinical adoption of machine learning models, predictive performance alone is not enough. Even a very accurate model has limited practical value, if its outputs cannot be understood or verified by its intended users. Clinicians need at least an explanation of which features drove a given prediction~\cite{carrasco_ribelles_2023_jamia}. For this reason the field of explainable AI (XAI) developed several post-hoc interpretation methods, for example SHAP (SHapley Additive exPlanations), LIME, or attention visualization. These methods show feature-level contributions to individual predictions in a form which a human reader can understand~\cite{moglia_2025_bmc}. For sequential models on administrative code sequences, interpreting predictions has additional problems. Recurrent architectures like LSTMs do not expose easily interpretable internal states. Transformer-based models give attention weights across the input sequence which can be visualized, but whether these weights faithfully correspond to model reasoning is still debated~\cite{nature_claims_transformer_2022}.

In the administrative healthcare setting, the target audience for machine learning outputs is usually not a data scientist. Healthcare professionals, epidemiologists, and real-world evidence specialists usually do not have the technical background to interpret raw prediction probabilities, confusion matrices, or log-loss curves. The gap between technically correct model evaluation and model communication which is usable by these audiences is one of the main barriers for practical deployment of AI in clinical workflows~\cite{carrasco_ribelles_2023_jamia, bmj_tripod_ai_2024}. So it is necessary to design interactive visual interfaces which translate model outputs into formats usable by non-technical end users. For example individual risk trajectories, population-level prediction distributions, or feature attribution plots for specific patients. The problems described in this section together show that a dedicated platform is needed. Such a platform should store and organize machine learning model outputs and also provide tools for accessible visualization and interpretation. The following sections describe existing tools that partially solve this problem, and then the general software engineering building blocks needed to implement a new one.


\subsection{Existing Tools and Platforms for EHR Data Analysis}
\label{sec:existing_tools}

The volume of longitudinal EHR data is growing. So several software tools for health data analysis were developed. Current tools are mostly split in two types: tools designed for technical research and tools designed for high-level clinical reporting. So there is a gap for a tool that would make machine learning outputs usable by clinical end users.

\subsubsection{Research-focused tools and environments}
Data scientists and machine learning engineers usually analyze EHR data with programmatic workflows. Python libraries like Pandas, Scikit-learn, and PyTorch are the core of these pipelines. They are often used in computational notebook environments like Jupyter~\cite{kluyver2016jupyter}. These environments are very flexible for model development and iterative data exploration. But they were designed as research and prototyping tools and not as production-grade clinical interfaces. Jupyter notebooks were originally made as a publishing format for reproducible computational workflows and not as deployable user-facing applications~\cite{kluyver2016jupyter}. They need the user to write or execute code, manage software dependencies, and configure computational environments. It is a too steep learning curve for non-technical healthcare professionals, epidemiologists, and real-world evidence specialists, who only need to interpret model outputs and not train them.

\subsubsection{Clinical analytics dashboards}
To connect raw data and end users, interactive web dashboards are used more and more. Frameworks like Dash (built on top of the Plotly visualization library) allow developers to make analytical web applications fully in Python. So data insights can be shown without users interacting with the underlying code~\cite{dash_plotly_framework}. Such frameworks were demonstrated in clinical settings through example applications which visualize patient volumes, waiting times, and care scores with interactive heatmaps and tables~\cite{dash_clinical_analytics_plotly}. These dashboards make descriptive statistics and operational health data more accessible. But they are usually hardcoded to specific static datasets and predefined metrics. They are not designed to dynamically accept, process, and visualize the complex, longitudinal outputs of multiple different machine learning models.

\subsubsection{Limitations and UX barriers in clinical decision support}
The main limitation of the current situation is that machine learning outputs are not connected to user-centered design (UCD). Recent systematic reviews of clinical decision support systems (CDSS) show that poor usability and inadequate user experience (UX) are major barriers for adoption of AI in healthcare~\cite{pubmed_40540451}. When machine learning outputs are presented, they often do not have contextual explanations. Or they are shown in formats that do not fit clinical workflows~\cite{jamia_bias_ui_2025}. To solve these barriers, the literature recommends practical UCD approaches which actively involve end users in the design process. So the interface can intuitively communicate model reliability, data quality, and feature importance~\cite{pubmed_40540451}. There is no tool that combines flexibility of ML model deployment with accessibility of user-centered clinical dashboards. It motivates the development of a dedicated data catalog application.

Another architectural aspect of deploying machine learning applications in clinical settings is differentiated access control. Healthcare web applications usually have users with different roles and authorization requirements. On one side there are system administrators or data scientists who configure and manage predictive models. On the other side there are clinical end-users, for example epidemiologists or real-world evidence specialists, who access and interpret model outputs. Role-Based Access Control (RBAC) is the established mechanism for managing such hierarchies in health information systems. In RBAC permissions are assigned to named roles and not to individual users. It enforces the principle of least privilege. Each user can access only the data and functionality needed for their function~\cite{kim2015rbac}. Implementation of RBAC in web-based healthcare applications has practical problems specific to the clinical domain. For example the design of role taxonomies that reflect organizational hierarchies, communication of access boundaries to end-users, and ongoing role maintenance when responsibilities change~\cite{samonte2024rbac}. These considerations directly influence the user and model management design of the application described in this work. Role-differentiated access determines which users can upload and configure models and which can only query and visualize prediction results.

\subsection{Web Applications for Machine Learning Model Deployment}
\label{sec:web_apps_ml}

To deploy a machine learning model outside a local notebook environment, the trained model artifacts have to be transformed into an accessible service. In modern software architectures it is usually done by serving the model over a network and providing a user-friendly frontend to interact with it.

\subsubsection{Model serving approaches}
The standard approach is to expose the model through a REST (Representational State Transfer) API. So client applications can send input data over HTTP and get predictions in real time. To make it easier, specialized model serving tools were developed. Platforms like MLflow provide lifecycle management. They allow models trained in different Python frameworks to be packaged in a standard format and deployed automatically as REST endpoints~\cite{mlflow_serving_2026}. Another option are open standards like ONNX (Open Neural Network Exchange). They allow models to be exported from the original training framework (for example PyTorch or TensorFlow) and executed in a highly optimized inference engine. So there is cross-platform compatibility and high performance in production environments~\cite{onnx_serving_2020}.

\paragraph*{Asynchronous task queues for long-running inference}

A practical problem when serving machine learning models through REST APIs is that inference for sequential models can take several seconds. It is especially true for models working over long patient histories. If such requests are served synchronously within an HTTP connection, connection timeouts can be exhausted. Application responsiveness is also degraded. A well-established pattern for decoupling long-running computation from synchronous HTTP responses is a distributed task queue. The HTTP server accepts the request, puts a prediction task into the queue, and immediately returns a task identifier to the client. A separate pool of worker processes consumes the queue and executes the inference. The client polls for the result later. This pattern is commonly implemented with Celery as the task execution framework and Redis as the message broker. So the queue is durable (tasks are not lost if one worker crashes) and it is easy to add more worker processes when the load grows~\cite{maksimovic2024taskqueue}.

Another important prerequisite for reproducible model deployment is environment isolation. We used Docker so the application runs in the same environment on any machine. With Docker the application and all its runtime dependencies (the Python interpreter, specific library versions, system libraries, and configuration files) are put together into one image. The image is then run as a container. The same container runs in development, during testing, and on the production server. So the deployment problems that happen when the server has different libraries or a different Python version than the developer machine are avoided~\cite{rivera2018containers}. This approach is also used in clinical informatics. The KETOS platform by Gruendner et al.\ uses Docker virtualization to give researchers reproducible data analysis and model development environments, accessible through web services. It supports the full research lifecycle, from data retrieval through model training to deployment as a clinical decision support service inside a hospital. So Docker-based workflows are viable for healthcare ML applications~\cite{gruendner2019ketos}. For multi-service applications which combine a web server, task queue, and several database backends, Docker Compose extends this pattern. All services and their interdependencies can be defined in one declarative configuration file. The whole application stack can be launched reproducibly with one command and without manual environment configuration~\cite{list2018dockercompose}.

\subsubsection{Backend frameworks for API development}
The flow of data between the database, model serving layer, and user interface has to be coordinated. For this purpose a reliable backend application is needed. In Python two frameworks are mostly used for this purpose: Flask and FastAPI. Flask is a mature, lightweight micro-framework. It is simple and flexible, so it is popular for straightforward web applications. But for machine learning workflows with heavy computational tasks or large data payloads, FastAPI is a better option~\cite{fastapi_vs_flask_2026}. FastAPI is built on ASGI (Asynchronous Server Gateway Interface). So it can handle multiple concurrent requests without blocking the execution thread. It is important for AI applications where model inference can take several seconds. FastAPI also provides automatic data validation and interactive API documentation out of the box. It reduces development overhead~\cite{fastapi_vs_flask_2026, geeks_fastapi_2024}.

Modern data-intensive web applications often use a polyglot persistence strategy. Instead of using one general-purpose database, different data management requirements inside one system are assigned to different storage technologies. Each storage technology fits one access pattern~\cite{khine2019polyglot}. In the healthcare application domain this strategy is well-established. Munoz et al.\ compared relational and NoSQL persistence mechanisms for standardized ISO/EN~13606 EHR data. They found that document-oriented NoSQL databases like MongoDB give better scalability and query performance for flexible, variable-structure records. These records are typical for clinical event sequences. Relational systems keep their advantage for well-defined, tabular data with complex integrity constraints. So it is useful to combine both in one system~\cite{munoz2017ehr}.

Another specific problem in healthcare web applications is efficient retrieval of relevant terms from large, hierarchically organized medical code vocabularies at interactive latency. Inverted-index engines like Elasticsearch are designed exactly for this access pattern. Wen et al.\ showed that context-aware indexing with Elasticsearch is much better than keyword-baseline approaches for retrieving relevant clinical concepts from large-scale EHR repositories. It uses the inverted index data structure, horizontal scaling, and information retrieval scoring algorithms to support fast and accurate concept search~\cite{wen2019elasticsearch}. Duan et al.\ implemented an Elasticsearch-backed hierarchical index over openEHR archetypes for patient screening in clinical research. They reported significant improvements in query performance and concept-matching accuracy over unindexed baseline approaches~\cite{duan2021openehr}. So the three-database architecture used in this work is justified. MongoDB is used for flexible storage of variable-length patient procedure code sequences. A relational database is used for structured management of users, models, and audit data. Elasticsearch is used for fast full-text search over the Czech administrative medical code vocabulary (\textit{Seznam zdravotních výkonů} and CZ-DRG codes).

\subsubsection{Frontend frameworks and interactive visualization}
On the client side, modern web applications use component-based JavaScript frameworks like React or Vue.js to build dynamic user interfaces. These frameworks manage the application state and update the browser view automatically, without full page reloads. So the user experience is smooth and app-like.

For healthcare analytics it is important to visualize complex, high-dimensional data. There are several JavaScript charting libraries which can be integrated with these frameworks. D3.js is the most powerful and flexible library. It allows developers to bind data directly to the Document Object Model (DOM) and create highly customized, interactive visualizations. But it has a steep learning curve. It can be complicated to integrate cleanly with React's own rendering engine~\cite{luzmo_react_charts_2025}. Higher-level libraries like Chart.js, Recharts, or Plotly.js provide pre-built, composable chart components. They give less granular control than D3.js, but they allow rapid implementation of responsive line charts, bar charts, and radar charts. These are needed for displaying patient trajectories and model confidence scores in clinical dashboards~\cite{luzmo_react_charts_2025}.

\subsubsection{Interactive Visualization of Sequential Model Outputs}
\label{subsec:interactive_viz_outputs}

Apart from general-purpose clinical dashboards, a smaller group of works deals specifically with the problem of presenting outputs of sequential probabilistic models in an interpretable, interactive form for clinical audiences. The DPVis system by Kwon et al.\ is an example. It integrates HMM training outputs into a coordinated set of interactive views. It shows state emission distributions, transition matrices, and per-patient state sequences. Clinical researchers can explore and compare disease progression pathways across patient subgroups~\cite{kwon2020dpvis}. DPVis was developed through a design study with clinicians, statisticians, and visualization experts. It shows that translation of probabilistic model outputs into accessible interfaces is feasible, but needs a lot of design effort.

For medical claims data, MedCV (Chanda et al., 2022) provides an interactive visualization environment for patient cohort identification. Clinical researchers can explore relationships between medical codes across patient timelines using word embedding-based similarity projections~\cite{chanda2022medcv}. MedCV focuses on cohort definition and not on predictive model outputs. But it shows that domain experts can meaningfully interact with complex, high-dimensional medical code data when they get appropriate visual abstractions.

These systems together show both the potential and the current limitations of interactive visualization for clinical data. They address specific, narrow tasks (progression analysis for one model, or cohort selection). They are implemented as standalone research tools and not as extensible web applications. They are also not designed to serve and compare outputs of multiple predictive models inside one unified interface. So this work aims to generalize the visualization of model outputs across multiple trained models in a deployable, role-appropriate web application.

\subsection{Data Catalog Concepts}
\label{sec:data_catalogs}

Systematic management of health data assets is more and more needed. One widely used answer to this need are the FAIR guiding principles. FAIR stands for findable, accessible, interoperable and reusable. These four properties were proposed by Wilkinson et al.\ in 2016 as a set of rules for how scientific data should be stored and described, so that other researchers can later find it and work with it~\cite{wilkinson2016fair}. In health informatics FAIR became the standard for data management. A scoping review of FAIR implementation practices across health data stewardship identified active FAIRification efforts in clinical registries, biobanks, hospital information systems and also national research infrastructures across almost all continents. Reported effects of these efforts include better data sharing between institutions and easier reuse of existing datasets for new studies~\cite{zeleke2022fair}. In the European context the FAIR principles are part of the H2020 Programme data management guidelines. They are also the basis of the European Open Science Cloud initiative. So they are a practical expectation for any health informatics project which handles research-grade data assets.

A data catalog is one of the main operational tools for implementing FAIR compliance in practice. It provides the metadata layer that the four FAIR properties actually rely on. Assets become findable because the catalog offers structured search over descriptions of the datasets. They become accessible because the catalog specifies how each asset can be obtained and under which permissions. They also become easier to interpret and reuse, because the catalog keeps the schema description and the provenance of each dataset next to the dataset itself. A recent taxonomy study of health metadata data catalogs (HMDCs) in European healthcare data environments confirmed that the design of such catalogs has to explicitly address FAIR principles. It is especially important for metadata asset design and data governance. The study identified dataset discoverability and security as the two main design problems for effective implementation~\cite{mallick2025hmdc}. This work develops a specialized instance of the catalog concept. It is focused not on raw EHR datasets, but on outputs of predictive models trained on longitudinal administrative sequences. It is a category of asset that existing general-purpose data catalogs do not address.

When organizations accumulate a lot of data across several systems, the challenge changes from data storage to data discoverability and governance. A data catalog is a software tool designed for this problem. It collects metadata (data about data) from the different storage systems that the organization uses and puts it into one place. So a user of the catalog can look up which datasets exist, what is inside them, where they came from, who owns them, and how they can be used, without having to query the underlying databases directly~\cite{data_world_catalog_2025}. The main goal of the catalog is therefore not to store data but to make it easier for people to get oriented in the data the organization already has~\cite{atlan_catalog_2025}.

\subsubsection{Existing data catalog platforms}
Several open-source and enterprise data catalog platforms are used today. For example, Amundsen was originally developed by Lyft. It focuses on automated discovery and search. It uses graph databases (like Neo4j or Apache Atlas) to map relationships between data assets and their users~\cite{automated_catalogs_2025}. DataHub was originally created by LinkedIn. It has a more complex architecture, built for production use in large companies and based on Kafka and Elasticsearch. It supports real-time metadata ingestion and lineage tracking~\cite{atlan_amundsen_datahub_2025}.

These platforms work well for traditional metadata management. They track schema changes, data lineage, and ownership of SQL tables. But they are designed for data engineers and data stewards. They focus on the technical structure of the data and not on what the individual data points actually mean clinically~\cite{matillion_catalog_2024}.

\subsubsection{The gap: Cataloging model outputs for clinical users}
Recently the concept of a ``machine learning data catalog'' appeared. It tries to capture the complex relationships between raw training data and deployed ML models~\cite{murdio_ml_catalog_2026}. But even these specialized catalogs are made for managing the ML lifecycle (for example tracking model versions and feature drift) for data scientists. They are not made for presenting actionable outputs to non-technical end users.

\subsection{Web Application Architectures and Frontend Principles}
\label{sec:web_arch_frontend}

Machine learning outputs become practically usable only when they are delivered through a user interface that fits the existing workflow of clinical end-users. The architectural and implementation choices of that interface determine responsiveness, maintainability, and deployability. Modern web applications are built out of established patterns that were developed over the last two decades. This section gives an overview of the main patterns and trade-offs that are relevant for this work: the division between traditional multi-page and single-page applications, the component-based frontend paradigm, virtual DOM and JSX as implementation techniques, rendering strategies (client-side, server-side, and hybrid), progressive web applications with service workers and browser storage, and the main state management paradigms.

\subsubsection{Single-page and multi-page applications}
\label{subsec:spa_mpa}

Traditional web applications were implemented as multi-page applications (MPAs). Each user action that required new content also required a full HTTP request-response cycle. The server re-rendered the whole page and returned new HTML. The browser then discarded the previous document. It is simple, server-driven, and friendly to search engines, because each URL resolves to a complete static HTML document that a crawler can index directly~\cite{olshevska2018spa_mpa, clean_commit_spa_mpa_2023}. MPAs also have small client-side footprint. So they work well on low-powered devices and unstable connections.

The single-page application (SPA) architecture appeared with the spread of asynchronous XMLHttpRequest calls (AJAX) and later with client-side frameworks like React, Angular, and Vue. In an SPA the browser loads one HTML shell and a JavaScript bundle once at the beginning. After that the application takes control of navigation. When the user moves between views, the application intercepts the navigation event, updates the browser URL via the History API, and re-renders only the components that changed. Data is exchanged with the backend as JSON over a REST or GraphQL interface~\cite{kornienko2021spa_secure, oneseventech_spa_mpa_2023}. So user interactions feel closer to a desktop application than to a document-oriented website. There is no full-page reload.

SPAs also have disadvantages. The initial JavaScript bundle is larger, so first-load time on slow networks can be worse than with an equivalent MPA. Search engine indexing is harder, because crawlers have to execute JavaScript to see the rendered content. Memory leaks are easier to introduce, because the same DOM context is kept alive across many user actions~\cite{olshevska2018spa_mpa}. A recent comparative study of multi-page and single-page architectures for internal business applications found that MPAs remain a more lightweight and simpler default for content-centric websites. SPAs are justified when the application is state-heavy or has long interactive sessions, for example multi-step forms, dashboards, or communication tools~\cite{clean_commit_spa_mpa_2023}. For clinical data catalogs and scoring workflows, which are exactly this kind of stateful, multi-step interaction, SPAs are the commonly preferred choice.

\subsubsection{Component-based development and the virtual DOM}
\label{subsec:component_vdom}

The dominant way of structuring modern SPAs is the component-based approach. The user interface is decomposed into small, reusable pieces called components. Each component encapsulates its own markup, styling, local state, and behaviour. Components are composed into a tree and communicate through properties (parent-to-child) and events or callbacks (child-to-parent)~\cite{jain2024cdd}. It replaces the older pattern of direct DOM manipulation through global scripts or jQuery selectors. So the code becomes modular, easier to test, and easier to reuse across projects. Component-driven development is now shared by all major frameworks (React, Vue, Angular) and also by the native HTML5 Web Components standard~\cite{jain2024cdd, shah2023webcomponents}.

In React specifically, components can be written as class components or as function components. Class components extend the \texttt{React.Component} base class and use lifecycle methods (\texttt{componentDidMount}, \texttt{componentDidUpdate}, \texttt{componentWillUnmount}) for side effects. They also use \texttt{this.state} and \texttt{this.setState} for local state. Function components are plain JavaScript functions that take properties and return JSX. Before React 16.8 they were stateless. Since React 16.8 the Hooks API was introduced. Function components can now manage state with \texttt{useState}, side effects with \texttt{useEffect}, and context with \texttt{useContext}. So the same functionality that needed class components before can be expressed in function components with less boilerplate~\cite{goli2019react_hooks}. Hooks also enable custom composition. Developers can extract reusable stateful logic into custom hooks and share it between components without higher-order component wrappers or render props. It is now the recommended style in modern React codebases~\cite{goli2019react_hooks, jain2024cdd}.

A frequently discussed implementation detail of React and Vue is the virtual DOM. Direct manipulation of the real Document Object Model is slow. Every write triggers layout, paint, and sometimes reflow of the whole subtree. The virtual DOM is an in-memory lightweight representation of the actual DOM tree. When component state changes, the framework builds a new virtual tree. It is compared with the previous virtual tree through a diffing algorithm. Only the computed minimal set of differences is then applied to the real DOM~\cite{bondarenko2025dom, bai2023million}. So the expensive real-DOM operations are batched and minimized. Angular uses a different approach called Incremental DOM. It skips the virtual tree and applies changes directly through a Change Detection mechanism. A recent comparative analysis found that the virtual DOM gives a simpler reactive programming model at the cost of extra memory and computation. Incremental DOM can reach higher peak performance with more deterministic update paths~\cite{bondarenko2025dom}. Compile-time optimizations like Svelte or Million.js remove the runtime virtual DOM completely and generate imperative update code directly, which reaches 100--300\% faster rendering in micro-benchmarks~\cite{bai2023million, putra2025frontend_comparison}. For typical data-oriented dashboards the virtual DOM abstraction is fast enough and the developer ergonomics are usually more important than the last millisecond of rendering time.

JSX is a syntax extension used mainly in React. It lets developers write markup that looks like HTML directly inside JavaScript. JSX is not interpreted by the browser. It is compiled at build time to standard JavaScript function calls (\texttt{React.createElement} or similar). So the component tree is fully expressible as data that can be analyzed, transformed, and diffed by the virtual DOM reconciler~\cite{goli2019react_hooks}. It gives two practical advantages. The markup is co-located with the component logic, so the full behavior of one UI piece is visible in one file. Static analysis tools like TypeScript can also type-check the JSX directly. So prop types, event handlers, and children are validated at compile time and not only at runtime.

\subsubsection{Rendering strategies: client-side, server-side, and hybrid}
\label{subsec:rendering_strategies}

An orthogonal architectural choice is where the HTML that the user sees is actually generated. In client-side rendering (CSR) the server returns an almost-empty HTML shell plus a JavaScript bundle. The browser then executes the JavaScript, fetches data, and builds the DOM locally. It is the default mode of SPAs. In server-side rendering (SSR) the server executes the same component tree, produces fully rendered HTML, and sends it to the browser. The page is usable immediately after the first HTML response. A separate JavaScript bundle then attaches interactivity to the already-rendered markup through a process called hydration~\cite{clapton2025csr_ssr, iskandar2020csr_ssr}. Static site generation (SSG) is a related variant in which the HTML is pre-rendered at build time and served as static files.

Controlled experiments show that the choice has measurable effects. SSR usually improves Time-to-Interactive (TTI) and Largest Contentful Paint (LCP), because the first meaningful HTML arrives sooner. It also improves search engine optimization (SEO), because crawlers see the complete document without executing JavaScript. Reported improvements for SSR over CSR are 58--62\% lower LCP and TTI on large payloads~\cite{clapton2025csr_ssr}. CSR can have better First Contentful Paint (FCP) for smaller datasets, because the initial HTML response is smaller and faster. Hybrid approaches like the ``islands architecture'' render most of the page as static HTML and selectively hydrate only interactive components. It keeps the SSR delivery benefits and reduces the JavaScript payload. Recent work on Modular Rendering with Adaptive Hydration (MRAH) in React applications reported further performance improvements through selective, priority-driven hydration of components~\cite{chen2025mrah}. For internal clinical applications, where SEO is not relevant and the audience is authenticated, pure CSR is usually enough. SSR becomes important for public-facing medical information sites and patient portals, where first-load performance and indexing matter.

\subsubsection{Progressive web applications, service workers, and browser storage}
\label{subsec:pwa_browser_storage}

A progressive web application (PWA) is a web application that uses a defined set of browser APIs to provide an experience close to a native mobile application. The main building blocks of a PWA are the web application manifest (a JSON file describing the name, icons, and theme colors), a service worker (a JavaScript file that runs on a separate thread and can intercept network requests), and HTTPS as a security prerequisite~\cite{muawwal2024pwa, popa2022pwa_health}. Service workers act as a client-side proxy between the application and the network. They can cache static assets, serve them offline, synchronize data in the background, and receive push notifications from the server~\cite{mdn_service_workers, webdev_service_workers}. Measurements of PWA implementations report average page load times 25--30\% faster with cache-enabled service workers, and full offline availability of the application shell~\cite{muawwal2024pwa}. PWAs are relevant for healthcare applications, because they reduce the development cost of maintaining separate native Android and iOS builds, and they can continue to work in low-connectivity hospital environments~\cite{popa2022pwa_health}.

On the client side, modern browsers give several mechanisms for storing data locally. Cookies are small key-value pairs (usually up to 4 KB) automatically sent with every HTTP request to the same origin. They are traditionally used for session identifiers and authentication tokens. Because they are sent on every request, they are also bandwidth-expensive and can leak through cross-site requests if not marked with appropriate security flags (\texttt{HttpOnly}, \texttt{Secure}, \texttt{SameSite}). The Web Storage API introduced \texttt{localStorage} and \texttt{sessionStorage}. Both are simple key-value stores scoped to one origin. They hold up to roughly 5--10 MB of string data per origin. The difference is lifetime. \texttt{localStorage} persists until explicitly cleared. \texttt{sessionStorage} is cleared when the browsing tab is closed. For larger structured data there is IndexedDB, a transactional object store that supports indexed queries and holds hundreds of megabytes to gigabytes, depending on the browser quota. The Cache Storage API, used with service workers, is a separate store specifically designed to cache HTTP request-response pairs for offline use~\cite{mdn_service_workers, webdev_service_workers}.

An additional, browser-managed cache layer works automatically. HTTP caching uses response headers like \texttt{Cache-Control}, \texttt{ETag}, and \texttt{Last-Modified} to decide whether a resource can be reused without a new network request. So static assets like JavaScript bundles, CSS files, and images are typically fetched once and then reused from cache across many sessions. These storage and caching mechanisms are complementary. Cookies are used for authentication tokens with short expiry. \texttt{sessionStorage} and \texttt{localStorage} for lightweight UI preferences or form drafts. IndexedDB for large client-side datasets. The HTTP cache for shared static assets. A well-designed application picks the right mechanism for each type of data so that both security and performance are respected.

\subsubsection{State management paradigms}
\label{subsec:state_management}

In a component tree, some state is purely local (for example the ``is this dropdown open'' flag inside one menu component). Other state has to be shared between components that are not in a direct parent-child relationship (for example the currently logged-in user, a theme selection, or cached server data). Passing shared state through props across many levels is called prop drilling. It makes the code brittle and hard to refactor. For this reason several state management paradigms were developed in the React ecosystem~\cite{goli2019react_hooks, putra2025api_react_laravel}.

The first and simplest approach is built into React itself. Local state is managed with the \texttt{useState} hook for primitive values and \texttt{useReducer} for more complex update logic, following the reducer pattern from functional programming. For shared state, the Context API provides a mechanism that lets a component tree expose values through a \texttt{Provider} and consume them with the \texttt{useContext} hook. So any descendant can read the value without explicit prop passing. The Context API is appropriate for relatively static global data like the authenticated user or the UI theme. It has an important performance limitation: every consumer of the context re-renders when the provided value changes, even if only one part of the value is actually relevant to that consumer~\cite{ameer_redux_context_2026, geeksforgeeks_state_mgmt_2025}.

For larger applications the dedicated state management library Redux became an industry reference. It centralizes the whole application state in one immutable store. State changes are described as plain action objects and applied through pure reducer functions. It enforces a strict unidirectional data flow and makes state transitions predictable, replayable, and debuggable with time-travel development tools~\cite{ameer_redux_context_2026}. The trade-off is boilerplate. Every piece of state needs actions, action creators, and reducer cases. Modern lightweight alternatives like Zustand, Jotai, and Recoil try to keep the predictability of Redux with less boilerplate. Atoms-and-selectors libraries like Recoil or Jotai split the state into small independently observable units. So a component only re-renders when the specific atom it reads actually changes~\cite{geeksforgeeks_state_mgmt_2025}. Another complementary approach is server-state libraries. For example TanStack Query (formerly React Query) or SWR specialize in caching, invalidating, and synchronizing remote data fetched from REST or GraphQL APIs. Server state is conceptually different from client-only UI state, because it has to be kept in sync with a remote source of truth~\cite{geeksforgeeks_state_mgmt_2025}.

A further pattern, which is especially useful for analytical applications, is URL-based state management. Filter values, pagination offsets, and selected entities are kept as query parameters in the browser URL. So the application state is shareable through a link, reproducible, and survives browser refresh. It is a robust replacement for ad-hoc in-memory state in data exploration dashboards, where users commonly need to bookmark or share a specific view of a dataset~\cite{putra2025api_react_laravel}. The choice between these paradigms depends on application size and data characteristics. For small applications local state and Context API are enough. For large applications with frequent cross-component updates a dedicated state library (Redux or a lightweight alternative) is used. For data-heavy, server-driven applications a combination of server-state caching and URL-based state often replaces the need for a global store~\cite{geeksforgeeks_state_mgmt_2025}.

\subsubsection{Browser compatibility and accessibility}
\label{subsec:compatibility}

A last frontend concern is that modern web applications run on many browser engines (Blink, WebKit, Gecko) and on many device classes, from mobile phones to workstations. Support for specific JavaScript features, CSS properties, and browser APIs varies across browser versions. Industry practice is to target a defined browserlist (for example ``last 2 versions of modern browsers''). A build toolchain like Babel or SWC then transpiles newer JavaScript features into backwards-compatible code for the chosen targets. PostCSS or similar tools add vendor prefixes to CSS. Polyfills fill missing Web API implementations. For healthcare applications that are used inside a controlled institutional environment the browser target can be narrowed. For public-facing medical websites it has to be much broader, because patients use a wide range of personal devices. Accessibility standards (Web Content Accessibility Guidelines, WCAG) are also relevant for healthcare, because users with visual or motor impairments must be able to interact with clinical tools. WCAG-conformant component libraries, semantic HTML, and correct ARIA attributes are the building blocks~\cite{yuan2024designrepair}.

\subsection{APIs and backend communication}
\label{sec:apis-backend}

A web application is usually split into two sides. One side is the browser code that the user sees, another side is the server code that holds the data and does the heavy work. These two sides talk to each other over the network through an API. So the design of the API decides how the whole system behaves. It controls what data can be asked for, how responses look, how errors are signaled, and how much work is done on each call. Because of this, the API style is one of the first decisions that has to be made when designing a web app \autocite{addanki2025api_shifts, marchuk2023db_access}.

There are several common API styles. The older one is SOAP, which is based on XML messages and strict contracts described by WSDL files. SOAP is still used in enterprise systems and in some healthcare integrations, but it is quite heavy and not so friendly for modern JavaScript clients. The most widely used style today is REST, which is lighter and works directly on top of HTTP. A newer alternative is GraphQL, where the client sends a query describing exactly which fields it wants and the server returns just that. GraphQL can reduce over-fetching and under-fetching of data, but it also moves a lot of complexity to the server. Another style is gRPC, which uses binary Protocol Buffers over HTTP/2 and is mostly used for service-to-service communication inside a cluster, not between a browser and a server \autocite{addanki2025api_shifts, sadaf2024rest_graphql, toman2020rest_soap, stepien2025rest_graphql, brito2020rest_graphql}.

\subsubsection{REST and Fielding constraints}
\label{sec:rest-constraints}

REST stands for Representational State Transfer. It was defined by Roy Fielding in his doctoral thesis and it is not a protocol, but an architectural style. REST uses HTTP as its transport and treats every piece of data as a resource with its own URL. So a resource like a patient record would have a URL such as \texttt{/patients/123}, and different HTTP methods are used to act on it. GET reads the resource, POST creates a new one, PUT or PATCH updates it, and DELETE removes it. The server returns the data in a chosen representation, usually JSON today, and the response carries an HTTP status code that says what happened \autocite{toman2020rest_soap, koziel2020rest_auth_security}.

Fielding defined six constraints that an API should follow to be called RESTful. The first is client-server separation, so that the user interface and the data storage can evolve independently. The second is statelessness, which means that each request carries all information needed to process it, and the server does not keep session state between calls. The third is cacheability, so that responses can say whether they can be cached and for how long. The fourth is a uniform interface, which gives every resource a stable URL and a fixed set of operations. The fifth is a layered system, where a client does not know if it talks directly to the server or through a proxy or a load balancer. The sixth constraint, code on demand, is optional and almost never used in practice \autocite{toman2020rest_soap, stepien2025rest_graphql}.

In practice, most REST APIs follow only some of these constraints. For example, true statelessness is hard when users have to log in, because some form of session has to exist somewhere. Also, not every API uses HTTP status codes consistently. So the term \enquote{RESTful} is often used loosely. Still, REST is the dominant style for public and internal web APIs, because it is simple to implement, easy to test with a browser or a tool like curl, and well supported by every framework \autocite{addanki2025api_shifts, toman2020rest_soap}.

\subsubsection{HTTP methods, status codes, and headers}
\label{sec:http-details}

Because REST sits on top of HTTP, understanding HTTP in detail is part of understanding web backends. HTTP is a request-response protocol, where the client sends a request with a method, a URL, headers, and an optional body, and the server sends back a status code, headers, and an optional body. The methods most used in APIs are GET for safe reads, POST for creating resources or for actions that have side effects, PUT for full replacement of a resource, PATCH for partial update, and DELETE for removal. HEAD and OPTIONS are also important, especially OPTIONS, which is used by browsers for CORS preflight checks \autocite{toman2020rest_soap}.

Status codes are grouped into five classes. Codes in the 2xx range mean success, for example 200 OK or 201 Created. The 3xx codes are for redirects. The 4xx codes say that the client made a mistake, like 400 Bad Request for malformed input, 401 Unauthorized when authentication is missing or wrong, 403 Forbidden when the user is logged in but not allowed, and 404 Not Found when the resource does not exist. The 5xx codes mean the server failed, usually 500 Internal Server Error or 503 Service Unavailable. A well-designed API uses these codes properly, because automatic clients and monitoring tools rely on them \autocite{koziel2020rest_auth_security}.

Headers carry metadata about the request or the response. Common request headers are \texttt{Authorization} for credentials, \texttt{Content-Type} to describe the body format, \texttt{Accept} to ask for a specific response format, and \texttt{Cookie} for session tokens. Response headers include \texttt{Set-Cookie} for writing cookies, \texttt{Cache-Control} for caching rules, and \texttt{Content-Length} and \texttt{Content-Type} for the body. Headers are also where CORS rules and security policies like HSTS or Content Security Policy are expressed \autocite{chinthalapudi2025security}.

\subsubsection{Middleware and the request-response pipeline}
\label{sec:middleware}

Modern backend frameworks organize request handling as a pipeline of middleware functions. A middleware is a small piece of code that sits between the raw HTTP request and the actual business logic. Each middleware can read the request, do something with it, and then pass it to the next middleware or stop the pipeline and return a response directly. The same happens in reverse for the response on the way back. So the whole request-response cycle is a chain of these small components, stacked in a specific order \autocite{kumar2022owasp_middleware}.

Middleware is used for many cross-cutting things. Typical examples are logging every request, parsing cookies and JSON bodies, checking authentication tokens, enforcing rate limits, adding CORS headers, compressing responses, handling exceptions, and measuring response time. The pattern is popular because it keeps these concerns out of the actual endpoint code. The business logic just gets a clean request and returns data, while the middleware layer takes care of everything around it. FastAPI, Express, Django, ASP.NET Core, and most other frameworks all use this same concept \autocite{kumar2022owasp_middleware}.

The order of middleware matters a lot. For example, an authentication middleware has to run before an authorization middleware, because we first need to know who the user is and only then decide what the user can do. Similarly, an exception handler has to wrap the whole pipeline, so that any error in later middleware is caught. CORS middleware usually runs very early, because the browser sends a preflight OPTIONS request that has to be answered before anything else. If the order is wrong, the whole app may look like it works but have subtle security or correctness bugs \autocite{kumar2022owasp_middleware, chinthalapudi2025security}.

\subsubsection{Same-origin policy and CORS}
\label{sec:cors}

Browsers enforce the same-origin policy. It means that a page served from one origin, defined by the combination of scheme, host, and port, cannot freely read data from another origin. This rule is one of the main defenses of the web, because without it any page could silently read the logged-in user's data from other sites. But when a frontend and a backend live on different origins, which is very common today, the app has to cross this boundary on purpose. This is done through Cross-Origin Resource Sharing, usually called CORS \autocite{cors_mdn_2026, cors_kong_2024}.

CORS is a set of HTTP headers that the server uses to tell the browser which origins are allowed to read its responses. For simple requests, the browser just adds an \texttt{Origin} header and the server answers with \texttt{Access-Control-Allow-Origin}. For requests that are not simple, for example those with custom headers or with a method like PUT or DELETE, the browser first sends a preflight OPTIONS request. The server has to answer this preflight with the allowed methods, allowed headers, and allowed origin, and only then the browser sends the real request. CORS is usually implemented as a middleware in the backend, because it needs to run on every request and also answer the OPTIONS preflight automatically \autocite{cors_mdn_2026, cors_kong_2024, chinthalapudi2025security}.

\subsubsection{Object-relational mapping}
\label{sec:orm}

Backend code often has to read and write rows in a relational database. Writing raw SQL everywhere works, but it is error-prone and couples the code to one specific database dialect. So most frameworks provide an object-relational mapper, usually called ORM. An ORM represents database tables as classes and rows as objects, and translates method calls into SQL queries. Python has SQLAlchemy and SQLModel, Java has Hibernate, JavaScript and TypeScript have Sequelize and TypeORM, and C\# has Entity Framework. These tools are widely used and cover most of the database access in modern web apps \autocite{hermanto2025orm_sql, ivanov2024orm_mysql, yan2016orm_wild}.

ORMs have clear advantages. Queries are written in the host programming language, so the IDE can check types and autocomplete field names. Migrations can be generated automatically when the model classes change. Different database engines can be supported by the same code with small configuration changes. ORMs also prevent a lot of SQL injection bugs, because values are always passed as bound parameters, not as string concatenation. So for typical CRUD code, ORMs reduce the amount of boilerplate and make the code more readable \autocite{ivanov2024orm_mysql, marchuk2023db_access}.

But ORMs also have downsides. The main one is the impedance mismatch between the object model in code and the relational model in the database. Some operations, especially joins and aggregations, are harder to express through an ORM than directly in SQL. Another well-known issue is the N+1 query problem, where reading a parent object and its children turns into one query for the parent and N extra queries for the children, instead of a single join. This kind of problem is easy to introduce and hard to notice until the app gets slow under load. Benchmarks report that raw SQL or lower-level query builders are often faster than ORMs for complex reads, while ORMs remain faster to develop for simple cases. So in practice, many projects use an ORM for most of the code and drop down to raw SQL for the few queries that really need it \autocite{hermanto2025orm_sql, ivanov2024orm_mysql, yan2016orm_wild}.

\subsubsection{Authentication and authorization}
\label{sec:auth}

Authentication answers the question \enquote{who is this user} and authorization answers \enquote{what is this user allowed to do}. These two concerns are separate, but they are almost always implemented together in the backend, usually as middleware that runs on every protected endpoint. Modern web apps use several standard mechanisms and it is common to combine them \autocite{pyroh2025auth, golla2026auth_mechanisms, chinthalapudi2025security}.

The oldest mechanism is session-based authentication. After the user logs in with a username and password, the server creates a session record on its side and sends the session ID back to the browser as a cookie. On every next request, the browser sends this cookie automatically, and the server looks up the session in memory or in a store like Redis. Sessions are simple and easy to revoke, because the server can just delete the session record. But they keep state on the server, which makes horizontal scaling harder, and they need extra protection against cross-site request forgery because cookies are sent automatically \autocite{pyroh2025auth, koziel2020rest_auth_security}.

A more recent mechanism is token-based authentication with JSON Web Tokens, known as JWT. A JWT is a signed piece of JSON that contains claims like the user ID, the roles, and an expiration time. After login, the server signs a JWT and returns it, and the client sends it back on every request, usually in the \texttt{Authorization} header with the \texttt{Bearer} scheme. The server can verify the signature without any database lookup, so JWT is called stateless. This makes it attractive for distributed APIs. But stateless also means that revoking a token before its expiration is complicated, because there is no server-side record to delete. So JWT-based systems often keep a short expiration time and use refresh tokens to renew sessions, or they add a separate revocation list \autocite{koziel2020rest_auth_security, chinthalapudi2025security, golla2026auth_mechanisms}.

For third-party login and for delegated access between services, the standard today is OAuth 2.0 together with OpenID Connect. OAuth 2.0 defines flows where a user gives an application limited access to their data on another service without sharing the password. The most used flow in web apps is the authorization code flow with PKCE. OpenID Connect adds an identity layer on top, where the authorization server also returns a signed ID token that tells the application who the user is. This is how \enquote{log in with Google} or \enquote{log in with Microsoft} works, and it is also how single sign-on is implemented inside organizations. OAuth has been formally analyzed and the main flows are considered secure when implemented correctly, but many real-world bugs come from small deviations from the specification \autocite{fett2016oauth_formal, golla2026auth_mechanisms, chinthalapudi2025security}.

For authorization itself, a very common pattern is role-based access control, where each user has one or more roles and each role has a set of permissions. The middleware then checks on every request whether the user's role includes the required permission. Other models like attribute-based access control exist, but in most web apps role-based access is enough \autocite{samonte2024rbac, kim2015rbac}.

\subsubsection{Password storage}
\label{sec:passwords}

Storing user passwords is a specific problem that deserves its own care. Passwords must never be stored in plain text. They also should not be stored with a simple hash like MD5 or SHA-1, because modern hardware can try billions of these hashes per second, which makes offline brute-force attacks realistic. So passwords are stored with a slow, memory-hard hash function designed for this purpose. Each password also gets a random salt, which is stored together with the hash and ensures that two users with the same password do not end up with the same hash \autocite{owasp_password_storage_2024}.

The three main functions used today are bcrypt, scrypt, and Argon2. Bcrypt is older and still widely used, with a configurable work factor that can be increased over time as hardware gets faster. Argon2 is newer and was selected in 2015 as the winner of the Password Hashing Competition. It has a memory-hard design that makes attacks with custom hardware much more expensive. OWASP currently recommends Argon2id as the first choice and bcrypt as an acceptable option for legacy systems. Some systems also add a pepper, which is an extra secret value mixed into the hash and stored separately from the database \autocite{owasp_password_storage_2024, phc_argon2_2015}.

\subsection{Data storage for web applications}
\label{sec:data-storage}

A web application has to persist many kinds of data. There are structured records like users, accounts, and permissions, where every row follows the same schema. There are semi-structured documents like patient sequences with variable length and optional fields. There are large binary blobs like images, CSV files, and serialized machine learning models. There is also transient data like session tokens, background-job state, and cached query results that do not need durability but have to be accessed very fast. No single storage system is good for all of these at the same time, so modern web apps usually combine several systems, each chosen for a specific role. This pattern is sometimes called polyglot persistence \autocite{khan2022sqlnosql, moniruzzaman2013nosql}.

\subsubsection{Relational databases}
\label{sec:relational-db}

Relational databases are still the default choice for structured data with well-known schemas. Data is stored in tables, each table has a fixed set of typed columns, and rows are linked through primary and foreign keys. The query language is SQL, which is standardized across engines. The main open-source engines used in web apps today are MySQL, PostgreSQL, and MariaDB, and on the commercial side Microsoft SQL Server and Oracle. Studies comparing these engines usually show small differences in single-query performance, with PostgreSQL doing better on complex analytical queries and MySQL being slightly faster for simple reads on small datasets \autocite{wodyk2020sql_frameworks, khan2022sqlnosql}.

The main strength of relational databases is the ACID guarantees: atomicity, consistency, isolation, and durability. A transaction either fully happens or does not happen at all, and concurrent transactions do not interfere with each other in ways that violate the schema. So for data that must always be correct, like account balances or user credentials, a relational database is the safe choice. Another strength is the expressiveness of SQL for joins and aggregations across several tables, which is how most reporting queries are written \autocite{khan2022sqlnosql, hassan2024big_data_db}.

The weakness of relational databases shows up with very large volumes of semi-structured data, or with highly variable schemas where the shape of a record can change per user. Horizontal scaling is also harder than in NoSQL systems, because sharding across servers has to preserve the consistency of joins and transactions. So many systems keep the core business data in a relational database and move big, flexible, or high-throughput data elsewhere \autocite{hassan2024big_data_db, nasution2023relational_nosql}.

\subsubsection{NoSQL databases and their types}
\label{sec:nosql-types}

NoSQL is an umbrella term for several families of databases that step away from the strict relational model. The literature and practice now recognize four main types, each with a different data model and different typical use cases \autocite{aguilar2023nosql_slr, gajiwala2024nosql_review, moniruzzaman2013nosql, asaad2020nosql}.

The first type is \emph{document stores}. Data is kept as JSON-like documents inside collections, and each document can have its own fields without a fixed schema. This is convenient when records have nested structures or vary a lot between instances. The most used document store is MongoDB. Document stores are a good fit for content like articles, user profiles, product catalogs, and event logs. They are also often used for semi-structured clinical data like patient sequences or EHR snapshots, because the shape of the data changes over time and adding a new field does not require a schema migration \autocite{aguilar2023nosql_slr, silalahi2018mongo_mysql, gyorodi2015nosql_web}.

The second type is \emph{key-value stores}. They behave like a giant hash map, where every value is retrieved by a unique key. There is no query language beyond get, put, and delete. Key-value stores are extremely fast because the access pattern is simple, so they are mostly used as caches, session stores, feature stores, or low-latency lookup tables. Redis and Memcached are the most common examples \autocite{srai2023redis_mda, li2018surfing_redis}.

The third type is \emph{column-family stores}, sometimes called wide-column stores. Data is stored by column rather than by row, which makes sequential scans over a single column very fast. This is a good fit for time-series data and for analytical workloads over huge datasets. Cassandra and HBase are the known systems in this family \autocite{gajiwala2024nosql_review, kumar2018nosql_traditional}.

The fourth type is \emph{graph databases}. Data is modeled as nodes and edges, and queries traverse these connections. This is useful for social networks, recommendation systems, fraud detection, and medical knowledge graphs. Neo4j is the dominant system, and alternatives like JanusGraph, TigerGraph, and Nebula Graph exist. Benchmarks on the LDBC Social Network Benchmark report that Neo4j has the best balance of query speed and memory use across typical workloads, but the field is young and the differences between engines are still large \autocite{monteiro2023graphdb, santos2024graphdb, besta2023graphdb}.

NoSQL systems usually relax some ACID properties in exchange for horizontal scalability. Many of them offer eventual consistency rather than strong consistency, so updates can take time to propagate across replicas. This is acceptable for a comment feed or a product view counter, but not for financial transactions. Some newer systems, sometimes called NewSQL, try to combine the scalability of NoSQL with the ACID guarantees of relational databases \autocite{patel2025newsql_multimodel, khan2022sqlnosql}.

\subsubsection{Full-text search engines}
\label{sec:search-engines}

A separate class of stores handles full-text search. Traditional databases are slow when asked to find rows by substring match, especially with fuzzy matching or ranking. So modern apps use dedicated search engines like Elasticsearch, OpenSearch, or Solr, all based on the Apache Lucene indexing library. These systems build inverted indices that map each word to the documents containing it, and they rank results using scoring functions like BM25 or TF-IDF. Search engines are often used as a secondary store: the primary data lives in a relational or document database, and a subset of it is indexed in the search engine for fast lookup. This is very common in medical applications where users need to search ICD codes, drug names, or procedure codes with typos and partial matches \autocite{wen2019elasticsearch}.

\subsubsection{In-memory stores and caching}
\label{sec:in-memory}

Disk access is orders of magnitude slower than RAM. So in-memory stores are used when very low latency matters. Redis is the most used in-memory store today. It supports several data structures beyond simple key-value pairs, like lists, sets, sorted sets, hashes, and streams. Typical use cases are caching query results, sharing user session data between stateless backend servers, rate limiting API calls, pub/sub messaging, and acting as the broker for background-job queues like Celery or RQ. Redis can also persist data to disk with snapshots or an append-only log, so it can be used as a primary database for some workloads, but most apps treat it as an ephemeral cache \autocite{redis_use_cases_2024, celery_redis_devto}.

The idea of caching applies at many layers of a web application. Browsers cache static assets following \texttt{Cache-Control} headers. Content delivery networks cache at the edge. Reverse proxies like Nginx or Varnish cache full HTTP responses. The backend can cache query results or computed views in Redis or Memcached. The database itself keeps a buffer pool in RAM. Each layer removes some work from the layer below, so a well-designed app spends most of its requests on cheap cache hits and reaches the database only when needed \autocite{redis_use_cases_2024}.

\subsubsection{Object storage and files}
\label{sec:object-storage}

Databases are not a good place for big binary files like images, videos, CSV dumps, trained machine learning models, or user uploads. Storing large blobs in a relational database inflates the row size, blocks the buffer pool, and makes backups slow. So the standard pattern is to keep binary files in a dedicated object store and keep only their metadata and URL in the database \autocite{alti2024nfs_s3, mondal2024object_big_data}.

The de facto object storage API is Amazon S3. Each object is identified by a bucket name and a key, and it can have custom metadata. Objects are typically retrieved over HTTP with signed URLs. Besides the AWS service itself, there are many S3-compatible implementations, notably MinIO, Ceph RADOS Gateway, and Google Cloud Storage and Azure Blob Storage through their own APIs. This S3 model is now used almost everywhere for static assets, data lakes, and ML artifacts \autocite{alti2024nfs_s3, gadban2021s3_hpc, mondal2024object_big_data}.

For legacy systems, traditional file storage options still exist: local disks, network file systems like NFS, and distributed file systems like Lustre or HDFS. Object stores have mostly replaced these for new cloud-native workloads, because they offer elastic scaling, high durability through replication, and pay-as-you-go pricing. Migration studies report that moving enterprise NFS workloads to S3 gives better scalability and lower cost, but requires adapting applications that assumed POSIX semantics, because S3 does not support append or partial writes in the same way as a file system \autocite{alti2024nfs_s3}.

\subsubsection{Storage of machine learning models}
\label{sec:ml-storage}

Storing and serving trained machine learning models is a specific subcategory of storage. A trained model is a file, so in principle it can be kept anywhere. In practice there are two problems. First, the file format matters for who can load the model. Scikit-learn models are typically saved as Python pickle files, which are fast and easy but only work back in Python and are known to be unsafe if loaded from untrusted sources because pickle can execute arbitrary code. PyTorch and TensorFlow have their own native formats. The Open Neural Network Exchange format, called ONNX, is a shared format for exporting models between frameworks and for serving them with language-independent runtimes \autocite{kubeflow_model_artifact_2024, fortra_pickle_2025}.

The second problem is versioning. Models are retrained periodically, different versions have different metrics, and it has to be clear which version is currently serving predictions. Ad-hoc solutions with filenames like \texttt{model\_v3\_final.pkl} do not scale. So dedicated model registries exist, most notably MLflow Model Registry. A registry assigns a unique name and version to every model, stores the actual artifact file in object storage, and keeps metadata like the training run, metrics, and stage (staging, production, archived). This makes it possible to roll back to a previous version, to compare versions, and to audit which model produced a given prediction \autocite{mlflow_registry_2025, lakefs_mlflow_2024}.

\subsection{Deployment and operations}
\label{sec:deployment}

Once a web application is built, it has to run somewhere. In the past, this meant installing the operating system, the runtime, the web server, and the application code directly on a physical or virtual machine. This approach works, but it makes every server slightly different, and a change that works on the developer's laptop may fail on the production server because of a different library version or a missing configuration file. This is the classic \enquote{it works on my machine} problem. Modern deployment practices try to remove this gap by packaging the application together with everything it needs to run, and by automating the path from source code to running service \autocite{prasetyo2021docker_web, galhotra2025containerization}.

\subsubsection{Containers and Docker}
\label{sec:containers}

A container is a lightweight form of virtualization. Unlike a virtual machine, a container does not include a full operating system. It shares the kernel with the host and packages only the application and its user-space dependencies. This makes containers much smaller and faster to start than VMs, often seconds instead of minutes, and allows many of them to run on the same machine without big overhead. Docker is the most used container engine, and it became the practical standard for packaging web applications \autocite{galhotra2025containerization, silva2024docker_lxd, rivera2018containers}.

The recipe for building a container image is written in a Dockerfile. It is a text file that describes the base image, the dependencies to install, the files to copy in, and the command to run. From the Dockerfile, the build process produces an image, which is an immutable layered filesystem. The same image can be run as many container instances as needed. Because the image contains the exact version of every dependency, the application behaves the same way on the developer's laptop, on a test server, and in production. This solves most of the environment-drift problems that used to plague manual deployments \autocite{prasetyo2021docker_web, galhotra2025containerization}.

A common pattern is the multi-stage build. The Dockerfile has two or more stages: a first stage with build tools, compilers, and dev dependencies that produces the compiled artifact (for example the built frontend bundle), and a final stage that only contains the runtime and the artifact. So the final image is small, does not ship build tools or source code, and has a smaller attack surface \autocite{galhotra2025containerization, list2018dockercompose}.

\subsubsection{Composing multiple services}
\label{sec:compose}

Real web applications almost never consist of a single container. There is usually a frontend, a backend, one or several databases, a cache like Redis, a task queue, and a reverse proxy in front of everything. Running each of these with a separate \texttt{docker run} command gets tedious very fast. So Docker Compose was introduced. It lets the whole topology be described in a single YAML file, with one block per service, the networks between them, the persistent volumes, and the environment variables. A single command then brings the whole stack up or down. Compose is well suited for development environments and for small single-host deployments \autocite{list2018dockercompose, galhotra2025containerization}.

For larger deployments, Docker Compose is not enough. It runs everything on one machine, does not handle failures, and does not distribute containers across a cluster. This is where container orchestration platforms come in.

\subsubsection{Kubernetes and orchestration}
\label{sec:kubernetes}

Kubernetes is the dominant container orchestration platform today. It was open-sourced by Google and is maintained by the Cloud Native Computing Foundation. Kubernetes takes a cluster of machines, called nodes, and treats them as a single pool of resources where containers can be scheduled. The user describes the desired state declaratively in YAML manifests: how many replicas of a service should run, how much CPU and memory each needs, how they should be exposed to the network, what secrets they can read. Kubernetes then works to make the real state match this description, restarting failed containers, rescheduling them to healthy nodes, rolling out new versions gradually, and rolling back on failure \autocite{lakka2025kubernetes, rodriguez2018orchestration, rao2022k8s_cicd}.

The building blocks are \emph{pods} (one or several containers that share a network namespace and always run together), \emph{deployments} (describing how many replicas of a pod should run and how updates are rolled out), \emph{services} (stable network endpoints in front of pods), and \emph{ingresses} (HTTP routing rules). Health probes tell Kubernetes when a container is alive and when it is ready to receive traffic, so broken instances are replaced automatically. Horizontal and vertical autoscalers adjust the number of replicas or the resources per replica based on observed load \autocite{lakka2025kubernetes, aldwyan2021elastic_k8s}.

Kubernetes is powerful, but it is also complex. Studies comparing it against simpler orchestrators like Docker Swarm report that Swarm is easier to set up and has lower overhead for small clusters, while Kubernetes scales further and has a much larger ecosystem of tools. For a small project, Docker Compose or Swarm is often enough; for a large distributed system with many services and strict availability requirements, Kubernetes is the usual choice \autocite{marella2024k8s_swarm, cochak2025k8s_docker}.

\subsubsection{Cloud platforms and serverless}
\label{sec:cloud-serverless}

Most modern web applications run on a cloud platform rather than on owned hardware. The three big public clouds are Amazon Web Services, Microsoft Azure, and Google Cloud Platform. They provide the usual building blocks: virtual machines, managed databases, object storage, load balancers, container registries, and managed Kubernetes clusters like EKS, AKS, and GKE. So a project can pick the level of abstraction it needs. At the lowest level, it rents VMs and installs everything itself. At the next level, it uses managed containers, where the cloud runs the cluster for it. At the highest level, it goes serverless \autocite{galhotra2025containerization, kotakonda2026k8s_ecs}.

Serverless computing, with services like AWS Lambda, Google Cloud Functions, and Azure Functions, abstracts the server away entirely. The developer writes a function, the cloud platform runs it on demand, scales it automatically, and bills per invocation. So the user only pays for actual execution time, and scaling to zero is possible when there is no traffic. This is attractive for event-driven workloads and for APIs with very bursty traffic \autocite{ponnaganti2025mlops, serverless_vs_containers_alibaba_2023}.

Serverless also has clear limitations. Cold start latency can be hundreds of milliseconds to several seconds when a function has not been called for a while. Execution time is usually capped, so long-running jobs do not fit. The runtime environment is restricted, and portability between providers is low because each cloud has its own specific APIs. Containers on the other hand have consistent performance once running, run anywhere Docker is available, and give full control of the environment, but they also require someone to manage scaling and failover. In practice many systems mix both: long-running stateful services run in containers, while bursty event-driven tasks go to serverless \autocite{serverless_vs_containers_alibaba_2023, serverless_vs_containers_cloudoptimo_2025, serverless_vs_containers_datadog_2024}.

\subsubsection{CI/CD pipelines}
\label{sec:cicd}

Deployment is not a one-off event. A live application is updated many times over its life, so the path from a new commit to a running production version has to be automated. This is done with Continuous Integration and Continuous Deployment pipelines, together called CI/CD \autocite{shahin2017cicd_slr, ali2023devops_cicd}.

Continuous Integration means that every code change is automatically built, its unit and integration tests are run, and the result is reported back to the developer within minutes. This catches regressions early, before they accumulate. Continuous Deployment then takes the green build and pushes it further: it builds the Docker image, pushes it to a registry, updates the running services in staging, runs end-to-end tests there, and finally rolls it out to production, usually with a gradual strategy like rolling update or canary release. Popular tools for this are GitHub Actions, GitLab CI, Jenkins, and CircleCI \autocite{singh2021cicd_cloud, istifarulah2023cicd_devops, goccia2022cicd_generator}.

A typical pipeline has several environment stages: development, where individual developers work, staging or test, which mirrors production as closely as possible, and production itself. Each stage has its own configuration, its own secrets, and often its own database. The promotion from one stage to the next is controlled by the pipeline, sometimes fully automatic and sometimes gated by manual approval. Infrastructure-as-code tools like Terraform or Pulumi are often used together with CI/CD, so that changes to infrastructure go through the same review-and-deploy flow as changes to application code \autocite{chatterjee2024cicd_devops, saxena2025devops_iac}.

\subsubsection{Reverse proxies and TLS}
\label{sec:reverse-proxy}

In almost every deployment there is a reverse proxy in front of the application servers. Nginx is the most common, but HAProxy, Caddy, Traefik, and Envoy are also widely used. The reverse proxy sits at the edge of the system and decides which backend gets each incoming request. So it is the place where many cross-cutting concerns are handled in one spot: TLS termination, HTTP/2 or HTTP/3 upgrade, request rate limiting, IP-based access control, caching of static assets, serving of the built frontend bundle, and routing between several backend services based on the URL path. It also hides the internal topology of the backend, because clients only ever see the proxy \autocite{nginx_docs_2025, devto_nginx_2024, solo_nginx}.

Modern web apps almost always use HTTPS. The reverse proxy terminates the TLS connection, decrypting the traffic and passing a plain HTTP request to the internal services. Certificates used to be expensive and manually renewed, but Let's Encrypt changed this by issuing free, short-lived certificates through the automated ACME protocol. Tools like Certbot, acme.sh, and Caddy's built-in automation obtain and renew certificates without human intervention. So enabling HTTPS on a new service is now a matter of a few minutes, and any app exposed to the internet is expected to have it. Additional hardening like HSTS, modern TLS ciphers, and HTTP Strict Transport Security is usually configured in the reverse proxy as well \autocite{lets_encrypt_automation_2024}.